% !TEX root = ../main.tex
% Partie 2 : on suit un pas de simulation de bout en bout. Chaque notion du
% cours apparait au moment ou on la rencontre. Le fil rouge (\fil{n}) en bas
% de chaque diapositive rappelle ou l'on est.

\section{Un pas de simulation}

% ============================================================ la carte
\begin{frame}{Le parcours d'un pas}
  \framesubtitle{huit étapes, de la scène à la souris}
  \centering
  \schema{parcours}
  \astcue{
    \cue{la \cuekey{carte} de toute la partie : on la suit étape par étape}
    \cue{1 à 3 : une seule fois \astsep{} 4 à 6 : deux fois par image \astsep{} 7 : une fois par image}
    \cue{pas fixe : \cuekey{Verlet suppose un dt constant}}
    \cuenext{étape 1 : d'où vient une scène}
  }
\end{frame}

% ============================================================ 1. scene
\begin{frame}{Une scène est une fonction}
  \framesubtitle{\code{std::function}, lambdas et un \code{template}}
  \centering
  \schema[5.3cm]{scene}
  \fil{1}
  \astcue{
    \cue{une scène = un nom, une taille et une \cuekey{fonction} qui remplit le monde}
    \cue{\code{R} recommence la scène : on rappelle \code{build} sur un monde neuf}
    \cue{les forces sont déclarées \cuekey{avant} \code{world} : détruites \cuekey{après} lui}
    \cuewarn{les particules gardent des références vers ces forces jusqu'au bout}
    \cuenext{ce que fait \code{build} : créer des particules}
  }
\end{frame}

% ============================================================ 2. particules
\begin{frame}{Créer une particule}
  \framesubtitle{constructeurs délégués et liste d'initialisation}
  \centering
  \schema[5.3cm]{particules}
  \fil{2}
  \astcue{
    \cue{trois constructeurs, \cuekey{un seul fait le travail} : les autres délèguent}
    \cue{\code{radius} est \code{const} : impossible de l'affecter dans le corps}
    \cue{\code{static} : le compteur appartient à la classe, pas à un objet}
    \cuenext{où vivent ces particules ? dans le \code{World}}
  }
\end{frame}

% ============================================================ 3. world
\begin{frame}{Le \code{World} possède tout}
  \framesubtitle{des adresses qui ne bougent jamais}
  \centering
  \schema[5.3cm]{world}
  \fil{3}
  \astcue{
    \cue{un ressort garde \cuekey{deux pointeurs} vers ses particules}
    \cue{\code{vector} : \code{push\_back} peut tout déplacer $\to$ pointeurs pendants}
    \cue{\code{deque} : les éléments ne bougent jamais}
    \cue{règle des trois : copie \cuekey{interdite} plutôt que copie profonde, choix justifié}
    \cuewarn{c'était le bug de la première version : ressorts sans effet}
    \cuenext{le monde est prêt, on lance un pas : les forces}
  }
\end{frame}

% ============================================================ 4. forces
\begin{frame}{Une particule somme ses forces}
  \framesubtitle{un seul appel, quatre comportements}
  \centering
  \schema[5.3cm]{forces}
  \fil{4}
  \astcue{
    \cue{chaque particule tient une liste de \code{Force\&} : \cuekey{liaison dynamique}}
    \cue{\code{Spring} s'inscrit lui-même sur ses deux particules}
    \cue{\code{World} ne fait aucun \code{if} sur le type d'une force}
    \cuewarn{limite : un ressort est calculé deux fois, une par extrémité}
    \cuenext{zoom sur le ressort : la formule s'écrit avec \code{Vec2}}
  }
\end{frame}

% ============================================================ 4b. Vec2
\begin{frame}[fragile]{Le ressort s'écrit comme la formule}
  \framesubtitle{surcharge d'opérateurs dans \code{Vec2}}
  \begin{columns}[c,onlytextwidth]
    \begin{column}{0.66\textwidth}
      \centering
      {\large
      $\vec F_A = \big[\,\textcolor{astblue}{k\,(\lVert\vec d\rVert - L_0)}
        + \textcolor{timaorange}{\rho\,(\vec v_B-\vec v_A)\cdot\hat d}\,\big]\,
        \textcolor{phelmared}{\hat d}$}
      \par\vspace{3mm}
\begin{lstlisting}[escapeinside={(*@}{@*)},basicstyle=\ttfamily\tiny]
const Vec2 d = ends[1]->getPos() - ends[0]->getPos();
const float length = d.norm();
const Vec2 (*@\textcolor{phelmared}{direction = d / length}@*);
const float (*@\textcolor{timaorange}{speed = (ends[1]->getSpeed() - ends[0]->getSpeed())}@*)
                    (*@\textcolor{timaorange}{.dot(direction)}@*);
return direction * ((*@\textcolor{astblue}{k * (length - L0)}@*) + (*@\textcolor{timaorange}{rho * speed}@*));
\end{lstlisting}
      \hfill\filechip{physics/forces/springs.cpp, abrégé}
    \end{column}
    \begin{column}{0.31\textwidth}
      \begin{tikzpicture}[x=1cm,y=1cm]
        \node[cls=datagreen/datafill,text width=3.9cm,align=left] (v) at (0,2.6)
          {\clsname{Vec2}\\[0.4ex]
           {\scriptsize méthodes :\\\code{+ - * / += -= == !=}\\[0.5ex]
            fonction libre :\\\code{float * Vec2}\\[0.5ex]
            \code{dot}, \code{norm}, \code{normalized}}};
        \node[astcallout,text width=3.7cm,font=\scriptsize] at (0,0.35)
          {\code{const Vec2\&} en paramètre\\méthodes \code{const} : lecture seule};
      \end{tikzpicture}
    \end{column}
  \end{columns}
  \fil{4}
  \astcue{
    \cue{chaque couleur de la formule se retrouve \cuekey{telle quelle} dans le code}
    \cue{\code{k * x} marche grâce à la fonction libre \code{operator*(float, Vec2)}}
    \cue{passage par \code{const\&} : pas de copie, pas de modification}
    \cuenext{on a l'accélération : il faut avancer dans le temps}
  }
\end{frame}

% ============================================================ 5. integration
\begin{frame}{Changer d'intégrateur en direct}
  \framesubtitle{\code{enum class}, une fabrique et une interface}
  \centering
  \schema[5.3cm]{integration}
  \fil{5}
  \astcue{
    \cue{\code{enum class} : un type à part, pas un \code{int} déguisé}
    \cue{toute la décision est dans la \cuekey{fabrique}, hors de \code{World}}
    \cue{démo : appuyer sur \cuekey{I} pendant le tissu}
    \cuenext{lequel est le meilleur ? on a mesuré l'énergie}
  }
\end{frame}

% ============================================================ 5b. energie
% [fragile] + \shorthandoff : pgfplots lit les << ; >> que babel-french rend actifs
\begin{frame}[fragile]{Euler ou Verlet : l'énergie}
  \framesubtitle{écart à l'énergie initiale, 10\,s, sans amortissement}
  \centering\shorthandoff{;:}%
  \begin{tikzpicture}
    \begin{axis}[astplot,name=a1,title={Pendule},xlabel={$t$ (s)},ylabel={écart (\%)},
        xmin=0,xmax=10,ymin=-1.2,ymax=1.2,
        legend pos=south west]
      \addplot[timaorange] table[x=t,y expr={(\thisrow{euler}/58.86-1)*100},col sep=comma]
        {figures/data/energie_pendule.csv};
      \addplot[astblue,dash pattern=on 3pt off 1.5pt] table[x=t,y expr={(\thisrow{verlet}/58.86-1)*100},col sep=comma]
        {figures/data/energie_pendule.csv};
      \legend{Euler,Verlet}
    \end{axis}
    \begin{axis}[astplot,at={(a1.east)},anchor=west,xshift=1.4cm,title={Double pendule (chaotique)},
        xlabel={$t$ (s)},xmin=0,xmax=10,ymin=-4,ymax=4]
      \addplot[timaorange] table[x=t,y expr={(\thisrow{euler}/100.847-1)*100},col sep=comma]
        {figures/data/energie_double_pendule.csv};
      \addplot[astblue,dash pattern=on 3pt off 1.5pt] table[x=t,y expr={(\thisrow{verlet}/100.847-1)*100},col sep=comma]
        {figures/data/energie_double_pendule.csv};
    \end{axis}
  \end{tikzpicture}
  \par\vspace{1mm}
  \begin{tikzpicture}
    \node[cls=datagreen/datafill,font=\scriptsize,align=center]
      {critère du sujet tenu : \textbf{moins de 5\,\%} sur 10\,s avec Verlet (pendule : 0,7\,\%)\\[0.2ex]
       {\tiny\color{textgrey}Euler \emph{semi-implicite} conserve aussi bien l'énergie : les deux schémas sont symplectiques}};
  \end{tikzpicture}
  \fil{5}
  \astcue{
    \cue{mesuré par le programme \code{energie} : $\frac12 mv^2 + mgh + \frac12 k(\ell-L_0)^2$}
    \cue{les deux courbes se superposent : \cuekey{pas de dérive}}
    \cuewarn{on attendait Euler moins bon : c'est vrai pour l'Euler \emph{explicite}, pas le semi-implicite}
    \cuenext{après l'intégration : on remet les particules dans la boîte}
  }
\end{frame}

% ============================================================ 6. contraintes
\begin{frame}{Les contraintes}
  \framesubtitle{bords du monde, rebond et points épinglés}
  \begin{columns}[c,onlytextwidth]
    \begin{column}{0.70\textwidth}
      \centering\schema[5.3cm]{contraintes}
    \end{column}
    \begin{column}{0.28\textwidth}
      \centering\scenecard{figures/scenes/scene_10.png}{Pluie de balles}{2.2cm}
    \end{column}
  \end{columns}
  \fil{6}
  \astcue{
    \cue{le monde est une \cuekey{boîte} : 4 bords au lieu d'un simple sol}
    \cue{coefficient de restitution $e$ : chaque balle a le sien}
    \cue{Verlet déduit la vitesse des positions : il faut corriger son historique}
    \cuenext{le pas est fini, on dessine}
  }
\end{frame}

% ============================================================ 7. dessin / diamant
\begin{frame}{Dessiner : le diamant}
  \framesubtitle{\code{Spring} est une \code{Force} et un \code{Drawable}}
  \centering
  \schema[5.3cm]{diamant}
  \fil{7}
  \astcue{
    \cue{\code{Force} et \code{Drawable} héritent \cuekey{virtuellement} d'\code{Entity}}
    \cue{mesure \code{sizeof} : 112 contre 144 octets, et un seul id}
    \cue{le moteur appelle \code{Renderer} : SFML reste une \cuekey{boîte noire}}
    \cuenext{dernière étape : l'utilisateur attrape une particule}
  }
\end{frame}

% ============================================================ 8. souris
\begin{frame}[fragile]{Attraper une particule à la souris}
  \framesubtitle{une lambda et un \code{dynamic\_cast}}
  \begin{columns}[c,onlytextwidth]
    \begin{column}{0.56\textwidth}
\begin{lstlisting}
Particle *World::pickParticle(const Vec2 &at,
                              float max_dist) {
  Particle *best = nullptr;
  float best_dist = max_dist;
  // la lambda capture best et best_dist
  auto closer = [&](Entity *e) {
    if (Particle *p = dynamic_cast<Particle *>(e)) {
      const float d = (p->getPos() - at).norm();
      if (d < best_dist) { best_dist = d; best = p; }
    }
  };
  std::vector<Entity *> all = entities();
  std::for_each(all.begin(), all.end(), closer);
  return best;
}
\end{lstlisting}
      \hfill\filechip{physics/world.cpp}
    \end{column}
    \begin{column}{0.40\textwidth}
      \centering
      \schema{souris}
    \end{column}
  \end{columns}
  \fil{8}
  \astcue{
    \cue{\code{dynamic\_cast} renvoie \code{nullptr} pour un ressort}
    \cue{la lambda capture par référence : \cuekey{[\&]}}
    \cue{relâcher = lancer : la vitesse de la souris est rendue à la particule}
    \cuenext{ce que tout ça donne : les scènes}
  }
\end{frame}
